\chapter{Introduction}
Charity advertisements are one of the primary lenses through which audiences in the Global North are exposed to narratives about the Global South. These narratives have frequently reproduced colonial rhetoric of the Global South’s inferiority and dependency. Within these advertisements, Africa has historically been over-represented. Consequently, scholarship on charity advertising has primarily centered on Africa, with a particular focus on image representations. The existing literature has demonstrated that charity advertising images promote negative ideas about the Global South by excluding examples of modernity, focusing on women and children as victims, and portraying the Global North as saviors and the Global South as helpless. However, the language accompanying images has been critically understudied, and existing studies are largely qualitative; this study addresses this gap by conducting a quantitative analysis of charity advertising language.

Charity advertising imagery of the Global South has long been dominated by images of extreme suffering, including illness and emaciation, often focusing on children. However, recent research suggests that this is changing following the Black Lives Matter Movement and a shift away from emotional humanitarian appeals to consumer-focused post-humanitarian advertising (Girling \& Adesina, 2024). Changing advertising strategies mitigated the use of pitiful imagery, but have been criticized for repackaging the white-savior narrative by focusing solely on the role of the donor and commercializing aid. These changing advertising practices necessitate updated research about colonial language within this new framework.

Studying the language of charity advertisements is significant because the images are not viewed in isolation. Instead, the language accompanying them serves as “anchorage,” shaping how the audience interprets the image and the advertisement’s overall message (Barthes, 1977). This means that charity advertisements do not need to show images of beneficiaries suffering; the accompanying text can suggest they are facing pain or hardship and perform a similar function. Whether language continues to propagate colonial ideologies in spite of image improvements is critical to study because charity advertisements that reinforce negative ideas about the Global South contribute to Global North notions that the Global South is characterized by archaic ideas and lifestyles, make some members of diaspora communities more reluctant to claim their African identities, and undermine long-term development goals.

This study hypothesizes that the text of charity advertisements is likely to reproduce stereotypes of the Global South, even as these stereotypes become less prevalent in visual representations. Postcolonial theory suggests that the Global North represents the “Other” (the Global South) through a lens of power dynamics that portray the Global North as more empowered and active and the Global South as more helpless and passive. The language of charity advertisements has received less scrutiny than the accompanying images and may reinforce colonial ideas through subtle word choices. Therefore, this study posits that, compared with advertisements featuring Global North beneficiaries, those featuring Global South beneficiaries have a higher proportion of vulnerability-related language. This study also hypothesizes that advertisements featuring Global South beneficiaries will use a higher proportion of active language when describing donor and charity actions than advertisements featuring Global North beneficiaries.

To achieve the goals of this research, the study quantitatively compares the language used in newspaper advertisements in the UK for Global North beneficiaries with that for Global South beneficiaries to determine whether there are statistically significant language differences in key themes associated with vulnerability and charity aid. This study focuses on advertisements distributed in the UK because it is a former colonizing power in the Global North with a large charity sector. This cross-sectional comparative design is suitable for this question because all charities must use persuasive language; using advertisements with Global North beneficiaries as a baseline can isolate whether thematic language density differs specifically for advertisements with Global South beneficiaries.
